\documentclass[../../../include/open-logic-section]{subfiles}

\begin{document}

\olfileid{sth}{z}{power}
\olsection{د فرعي سټونو سټونه}

د يو بل بديهي اصل په وړاندې کولو سره مخته ځو:

\begin{axiom}[د فرعي سټونو سټونه]
د هر سټ $A$ دپاره سټ $\Pow{A} = \Setabs{x}{x \subseteq A}$ شتون لري.
\[
	\forall A \exists P \forall x(x \in P \liff (\forall z \in x)z \in A)
\]
\end{axiom}

د دې اصل توجيه په نسبتاً نېغه توګه کېږي. فرض کړئ چې $A$ په
پړاو $S$ کښې جوړ شوے دے. نو د $A$ ټول غړي له $S$ نه مخکې
موجود وو (د \stagesacc له مخې). په دې توګه، د بېلونې په توجيه
کښې زموږ د استدلال په شان، د $A$ هر فرعي سټ تر پړاو $S$ پورې
جوړېږي. نو دا ټول له $S$ نه په هر وروسته پړاو کښې موجود دي،
څو په يوه سټ کښې راټول شي. او پوهېږو چې داسې يو پړاو شته،
ځکه چې $S$ وروستے پړاو نۀ دے (د \stagessucc له مخې).
نو $\Pow{A}$ شتون لري.

دلته د فرعي سټونو د سټونو يوه ښکلې پايله ده:

\begin{prop}\label{thm:Products}
د هر دوو راکړل شويو سټونو $A, B$ دپاره د هغو کارتيزي حاصل ضرب $A \times B$ شتون لري.
\end{prop}

\begin{proof}
سټ $\Pow{\Pow{A \cup B}}$ د فرعي سټونو د سټونو او
\olref[sth][z][pairs]{prop:pairsconsequences} له مخې شتون لري.
نو د بېلونې له مخې دا سټ شتون لري:
\[
	C = \Setabs{z \in \Pow{\Pow{A \cup B}}}{(\exists x \in A)(\exists y \in B) z = \tuple{x, y}}.
\]
اوس د هر $x \in A$ او $y \in B$ دپاره سټ $\tuple{x, y}$ د
\olref[sth][z][pairs]{prop:pairsconsequences} له مخې شتون لري.
پر دې سربېره، ځکه چې $x, y
\in A \cup B$ دي، نو $\{x\}, \{x, y\} \in \Pow{A \cup B}$ او
$\tuple{x,y} \in \Pow{\Pow{A \cup B}}$ لرو. نو $A \times B = C$.
\end{proof}

په دې ثبوت کښې د فرعي سټونو اصل له بېلونې سره يوځاے کار کوي.
او دا د حيرانتيا خبره نۀ ده. بې له بېلونې به د فرعي سټونو اصل
دومره \emph{پياوړے} اصل نۀ و. ځکه چې بېلونه راته وايي چې د يوه
سټ کوم فرعي سټونه شتون لري، او په دې توګه ټاکي چې د فرعي سټونو
هر سټ څومره «غټ» دے.

\begin{prob}
وښيئ چې د هر دوو سټونو $A, B$ دپاره: (i) د هغو ټولو اړيکو سټ
شتون لري چې د تعريف ساحه ئې $A$ او د قيمتونو سټ ئې $B$ وي؛ او
(ii) له $A$ څخه $B$ ته د ټولو تابعو سټ شتون لري.
\end{prob}

\begin{prob}
$A$ يو سټ او $\sim$ په $A$ باندې د هم‌ارزۍ اړيکه وګرځوئ.
ثابت کړئ چې د $\sim$ له مخې په $A$ باندې د هم‌ارزۍ د ټولګيو
سټ، يعنې $\equivclass{A}{\sim}$، شتون لري.
\end{prob}

\end{document}
